\chapter{Vaccination}

\section*{Overview}
Vaccination involves administering a vaccine (antigenic material) to stimulate an individual's immune system to develop adaptive immunity to a pathogen, preventing infectious disease.

\section*{Alternative Names}
Immunization, Vaccine injection

\section*{Principle}
A vaccine introduces a killed or weakened organism, a subunit, or genetic material coding for an antigen without causing disease. The immune system mounts a primary response and produces memory cells, so later exposure to the real pathogen triggers a rapid, protective secondary response.
\section*{History}
Edward Jenner developed the smallpox vaccine in 1796. Louis Pasteur developed vaccines for rabies and anthrax. Jonas Salk developed the polio vaccine in 1952.

\section*{Why It Is Done}
Ordered to prevent infectious diseases (like influenza, tetanus, measles, COVID-19), protect public health, and satisfy travel or employment requirements.

\section*{Is This a Primary or Secondary Test or Procedure}
Primary preventive intervention for infectious diseases.

\section*{How It Is Performed}
Usually injected into a muscle (e.g., deltoid or anterolateral thigh) or given orally/nasally. The injection site is cleaned with alcohol beforehand.

\section*{Preparation}
Verify the patient's vaccination record and check for allergies (e.g., egg allergy for certain flu vaccines).

\section*{Contraindications and Precautions}
Severe allergic reaction (anaphylaxis) to a prior dose of the same vaccine. Live vaccines are contraindicated in pregnancy and immunodeficiency.

\section*{Risks}
Common side effects include pain/redness at the injection site, mild fever, or muscle aches. Anaphylaxis is extremely rare (~1 in 1 million).

\section*{Aftercare and Recovery}
Monitor the patient for 15 minutes after the injection for any immediate allergic reaction or vasovagal syncope.

\section*{Results and Interpretation}
Stimulates the production of antibodies, providing immunity. Some vaccines require checking antibody titers later.

\section*{Expected Results}
Successful administration of the vaccine; expected mild localized immune response without systemic complications.

\section*{Follow-up}
Documentation in immunization registry; schedule subsequent booster doses.

\section*{References}
\begin{enumerate}
  \item MedlinePlus. Vaccination. U.S. National Library of Medicine; 2025.
  \item Current Medical Diagnosis and Treatment. McGraw-Hill; 2025.
\end{enumerate}

\clearpage
